\DocumentMetadata{
  lang=en
} 
\documentclass[12pt]{report} 
 
% -------------------------------------------------- 
% Accessibility and document metadata 
% -------------------------------------------------- 

\usepackage[english]{babel} 
\usepackage[T1]{fontenc} 
\usepackage[utf8]{inputenc} 
 
% -------------------------------------------------- 
% Page layout and spacing 
% -------------------------------------------------- 
\usepackage[margin=1in]{geometry} 
\usepackage{setspace} 
\doublespacing 
 
% -------------------------------------------------- 
% Useful packages 
% -------------------------------------------------- 
\usepackage{graphicx} 
\usepackage{booktabs} 
\usepackage{amsmath, amssymb} 
\usepackage{hyperref} 
\usepackage{caption} 
\usepackage{float} 
 
% -------------------------------------------------- 
% Hyperref setup 
% -------------------------------------------------- 
\hypersetup{ 
    colorlinks=true, 
    linkcolor=black, 
    urlcolor=blue, 
    citecolor=black, 
    pdftitle={Accessible Thesis Example}, 
    pdfauthor={Your Full Name}, 
    pdfsubject={Sample Accessible Thesis in LaTeX}, 
    pdfkeywords={LaTeX, thesis, accessibility, WCAG} 
} 
 
\begin{document} 
 
% -------------------------------------------------- 
% Title Page 
% -------------------------------------------------- 
\begin{titlepage} 
    \centering 
    {\Large THE UNIVERSITY OF TEXAS AT ARLINGTON\par} 
    \vspace{1.5cm} 
     
    {\Huge\bfseries Accessible Thesis Example in LaTeX\par} 
    \vspace{1.5cm} 
     
    {\Large Your Full Legal Name\par} 
    \vspace{1.5cm} 
     
    {\large Thesis\par} 
    \vspace{1cm} 
     
    {\large Submitted in partial fulfillment of the requirements\\ 
    for the degree of Master of Science in Computer Science\par} 
    \vspace{1.5cm} 
     
    {\large May 2026\par} 
    \vspace{1cm} 
     
    {\large Arlington, Texas\par} 
\end{titlepage} 
 
% -------------------------------------------------- 
% Front Matter 
% -------------------------------------------------- 
\pagenumbering{roman} 
 
\chapter*{Abstract} 
\addcontentsline{toc}{chapter}{Abstract} 
 
This sample thesis demonstrates how a student can structure a dissertation or thesis in LaTeX while keeping accessibility in mind. The document uses semantic sectioning commands such as \verb|\chapter| and \verb|\section|, includes properly labeled tables and figures, and uses standard mathematical environments for formulas. The goal of this example is to provide a student-friendly model that supports both academic formatting and accessibility best practices. By following this structure, students can create documents that are easier to navigate, more readable, and more compatible with assistive technologies. 
 
\chapter*{Acknowledgements} 
\addcontentsline{toc}{chapter}{Acknowledgements} 
 
I would like to thank my advisor, committee members, and colleagues for their guidance and support throughout this work. 
 
\tableofcontents 
\listoftables 
\listoffigures 
 
% -------------------------------------------------- 
% Main Matter 
% -------------------------------------------------- 
\clearpage 
\pagenumbering{arabic} 
 
\chapter{Introduction} 
\section{Background} 
 
Accessibility in academic documents is important because theses and dissertations are often shared digitally as PDF files. If the structure is not created correctly, the document may look visually correct while still being difficult for screen readers and keyboard-only users to navigate. 
 
\section{Purpose of This Example} 
 
This thesis serves as a model for students who are learning to write in LaTeX. It demonstrates proper use of chapters, sections, tables, figures, equations, and references. 
 
\section{Organization of the Thesis} 
 
This thesis is organized into five chapters. Chapter 1 introduces the topic. Chapter 2 reviews prior work. Chapter 3 explains the methodology. Chapter 4 presents results, including tables, figures, and equations. Chapter 5 concludes the thesis. 
 
\chapter{Literature Review} 
\section{Accessibility in Digital Documents} 
 
Digital accessibility focuses on making content usable by all readers, including people who use assistive technologies. According to the Web Content Accessibility Guidelines, information should be perceivable, operable, understandable, and robust \cite{wcag21}. 
 
\section{LaTeX and Accessibility} 
 
LaTeX provides strong support for structured documents. Commands such as \verb|\chapter|, \verb|\section|, and \verb|\caption| give the document semantic organization. This is one reason LaTeX can be a strong tool for accessible academic writing when used carefully. 
 
\chapter{Methodology} 
\section{Document Structure} 
 
The first step in creating an accessible thesis is using the correct LaTeX structure. Students should avoid manually bolding text for headings and instead use sectioning commands. 
 
\section{Tables} 
 
Tables should be simple, include a caption, and include a label so they can be referenced clearly in the text. 
 
\begin{table}[H] 
\centering 
\caption{Average Reading Time by Document Type} 
\label{tab:readingtime} 
\begin{tabular}{lc} 
\toprule 
Document Type & Average Time (minutes) \\ 
\midrule 
Word Document & 18 \\ 
PDF Document  & 22 \\ 
Web Page      & 15 \\ 
\bottomrule 
\end{tabular} 
\end{table} 
 
Table~\ref{tab:readingtime} shows a simple example of a structured table with a caption and label. 
 
\section{Figures} 
 
Figures should include meaningful captions and should be referred to in the text. 
 
\begin{figure}[H] 
\centering 
\includegraphics[width=0.6\textwidth]{example_figure.png} 
\caption{Example figure showing a simple accessibility workflow} 
\label{fig:workflow} 
\end{figure} 
 
Figure~\ref{fig:workflow} illustrates how an example image should be inserted and labeled. 
 
\section{Mathematical Formulas} 
 
Mathematical expressions should be written using LaTeX math environments instead of inserting screenshots of formulas. 
 
An inline equation can appear in a sentence, such as $E = mc^2$. 
 
A displayed equation can be written as: 
 
\begin{equation} 
E = mc^2 
\label{eq:energy} 
\end{equation} 
 
Equation~\ref{eq:energy} is an example of a numbered equation that can be referenced in the text. 
 
A multi-line equation can be written using the \texttt{align} environment: 
 
\begin{align} 
f(x) &= x^2 + 2x + 1 \\ 
     &= (x+1)^2 
\end{align} 
 
\chapter{Results and Discussion} 
\section{Sample Findings} 
 
This chapter demonstrates how students can discuss tables, figures, and equations together in a real thesis chapter. 
 
The data in Table~\ref{tab:readingtime} suggests that reading time varies by document type. Figure~\ref{fig:workflow} provides a visual example of how accessibility review may proceed. Equation~\ref{eq:energy} demonstrates that formulas can be written in a searchable and scalable format. 
 
\section{Reference Management} 
 
Students should avoid manually typing citations and references. Instead, they should use BibTeX or BibLaTeX. For example, accessibility standards are cited in this thesis using the \verb|\cite| command \cite{section508}. 
 
\chapter{Conclusion} 
\section{Summary} 
 
This sample thesis demonstrates a beginner-friendly format for writing an accessible thesis in LaTeX. It includes front matter, structured chapters, a table, a figure, equations, and references. 
 
\section{Recommendations} 
 
Students should always use semantic section commands, include captions and labels, use descriptive references, and rely on standard LaTeX environments instead of manual formatting. 
 
% -------------------------------------------------- 
% References 
% -------------------------------------------------- 
\bibliographystyle{plain} 
\bibliography{references} 
 
\end{document} 
